\uuid{hvCC}
\titre{Notion d'équation différentielle}
\auteur{bodin ; croizat ; sacre}
\organisation{exo7}
\licence{CC BY-NC-SA 4.0}
\datecreate{2022-03-08}
\niveau{L1}
\module{Analyse}
\chapitre{Equation différentielle}
\sousChapitre{Résolution d'équation différentielle du premier ordre}
\difficulte{3}
\theme{Notion d'équation différentielle}
\exo7qcmid{hilisit.equations-differentielles.022}

\contenu{
\begin{question}
\qclassification{225.01}
Soit \(f\) une solution de l'équation différentielle \(y'=y+3x\). On sait de plus que la courbe représentative de \(f\) passe par le point \(A(-1,2)\). Quelles sont les affirmations exactes ?
\begin{answers}
    \good{La pente de la tangente à la courbe de \(f\) au point \(A\) est \(-1\).}
	\bad{La pente de la tangente à la courbe de \(f\) au point \(A\) est \(4\).}
	\good{La tangente à la courbe de \(f\) au point \(A\) admet pour équation : \(y=-x+1\).}
	\bad{La tangente à la courbe de \(f\) au point \(A\) admet pour équation : \(y=4x+6\).}
\end{answers}
\begin{explanations}
 La pente de la tangente au point $A(-1,2)$ est le nombre $f'(-1)$. Or on sait que \(f(-1)=2\) puisque la courbe représentative de \(f\) passe par \(A(-1,2)\). De plus, comme \(f\) est solution de l'équation différentielle \(y'=y+3x\), en considérant cette égalité pour la fonction $f$ et pour $x=-1$, on a : \(f'(-1)=f(-1)+3\times (-1)=2-3=-1\). La pente de la tangente en \(A\) est donc \(-1\). Enfin, les coordonnées du point \(A\) vérifient l'équation de cette tangente, ce qui permet d'obtenir que l'ordonnée à l'origine vaut bien $+1$ (on sait aussi  plus directement que l'équation de la tangente est $y = (-1) (x-(-1))+1 = -x+1$).
\end{explanations}
\end{question}
}
